% Compile with: uplatex file.tex && dvipdfmx file.dvi
%           or: latexmk -uplatex file.tex
\documentclass[dvipdfmx]{article}
\usepackage{amsmath, amssymb}
\usepackage{tikz}

\title{scenario\_poisson\_bracket}
\date{\today}

\begin{document}
\maketitle

To compute a Poisson bracket, the system must first know which coordinate and momentum are canonically conjugate. Declare $p_r$ as the momentum conjugate to the radial coordinate $r$.

Define $p_r$ as the canonical momentum of r

Registered canonical pair (r, p\_r). It is now used by "compute the Poisson bracket of f and g".

Declare $p_\theta$ as the momentum conjugate to the angular coordinate $\theta$.

Define $p_{\theta}$ as the canonical momentum of theta

Registered canonical pair (theta, p\_theta). It is now used by "compute the Poisson bracket of f and g".

Compute the basic Poisson bracket $\{r,p_r\}$. Since these are a canonical pair, the canonical commutation relation $\{q,p\}=1$ means this should equal 1.

Find the Poisson bracket of r and $p_r$

\begin{equation}
  1
  \label{eq:E1}
\end{equation}

Define the central-force Hamiltonian $H=p_r^2/(2m)+p_\theta^2/(2mr^2)+U(r)$ (the sum of the radial and angular kinetic energies and the potential $U(r)$).

$H = p_r^2/(2 m) + p_{\theta}^2/(2 m r^2) + U(r)$

\begin{equation}
  H = p_r^2/(2 m) + p_{\theta}^2/(2 m r^2) + U(r)
  \label{eq:E2}
\end{equation}

Define the angular momentum $L_z=p_\theta$ (which is simply the momentum conjugate to the angle).

$L_z = p_{\theta}$

\begin{equation}
  L_z = p_{\theta}
  \label{eq:E3}
\end{equation}

Compute the Poisson bracket of the Hamiltonian $H$ and the angular momentum $L_z$. Since $H$ does not depend on the angle $\theta$ (rotational symmetry), this vanishing confirms that $L_z$ is conserved.

Find the Poisson bracket of H and $L_z$

\begin{equation}
  0
  \label{eq:E4}
\end{equation}

Using Hamilton's equation of motion $\dot q=\{q,H\}$, compute the Poisson bracket of $H$ and the radial coordinate $r$ to find its time evolution.

Find the Poisson bracket of H and r

\begin{equation}
  -\frac{p_{r}}{m}
  \label{eq:E5}
\end{equation}

\end{document}